\ifdefined\gkver\else\def\gkver{student}\fi
\documentclass[\gkver]{gaokaozhenti}
\newcommand{\ccnum}[1]{{\fontspec{Noto Serif CJK SC}#1}}
\begin{document}

\chapter{2011年山东理综}
%% paperType: 理综物理部分

\begin{enumerate}
\item
%% number: 16
%% paperName: 2011年山东理综第 16 题 4 分
%% typeId: 2
%% type: 多选题
%% chapter: 其他
%% point: 物理学史
%% method: 无
%% score: 4
%% degree: 650
%% duplicateId: 0
%% body:
了解物理规律的发现过程，学会像科学家那样观察和思考，往往比掌握知识本身更重要。以下符合事实的是\xzanswer{AB}
\fourchoices[answer=AB]
{焦耳发现了电流热效应的规律}
{库仑总结出了点电荷间相互作用的规律}
{楞次发现了电流的磁效应，拉开了研究电与磁相互关系的序幕}
{牛顿将斜面实验的结论合理外推，间接证明了自由落体运动是匀变速直线运动}

%% answer: AB
%% memo:
\memoanswer{
1840 年英国科学家焦耳发现了电流热效应的规律；库仑总结出了点电荷间相互作用的规律；法拉第发现了电磁感应现象，拉开了研究电与磁关系的序幕；伽利略通过将斜面实验合理外推，间接证明了自由落体运动的规律。
}

\item
%% number: 17
%% paperName: 2011年山东理综第 17 题 4 分
%% typeId: 2
%% type: 多选题
%% chapter: 曲线运动
%% point: 人造地球卫星
%% method: 无
%% score: 4
%% degree: 650
%% duplicateId: 0
%% body:
甲、乙为两颗地球卫星，其中甲为地球同步卫星，乙的运行高度低于甲的运行高度，两卫星轨道均可视为圆轨道。以下判断正确的是\xzanswer{AC}
\fourchoices[answer=AC]
{甲的周期大于乙的周期}
{乙的速度大于第一宇宙速度}
{甲的加速度小于乙的加速度}
{甲在运行时能经过北极的正上方}

%% answer: AC
%% memo:
\memoanswer{
对地球卫星，万有引力提供其做圆周运动的向心力，则有 $G\frac{Mm}{r^{2}}=m\frac{v^{2}}{r}=mr\omega^{2}=mr\frac{4\pi^{2}}{T^{2}}=ma_{n}$，可知半径越大速度越小，半径越大加速度越小，同步卫星的轨道与赤道共面，第一宇宙速度为最大环绕速度，可见 A 正确、C 正确。
}

\item
%% number: 18
%% paperName: 2011年山东理综第 18 题 4 分
%% typeId: 1
%% type: 单选题
%% chapter: 直线运动
%% point: 竖直上下抛运动
%% method: 无
%% score: 4
%% degree: 450
%% duplicateId: 0
%% body:
如图所示，将小球a从地面以初速度$v_{0}$竖直上抛的同时，将另一相同质量的小球b从距地面$h$处由静止释放，两球恰在$\frac{h}{2}$处相遇（不计空气阻力）。则\xzanswer{C}

\onepicture[w=1.80cm]{figs/18.png}

\fourchoices[answer=C]
{两球同时落地}
{相遇时两球速度大小相等}
{从开始运动到相遇，球a动能的减少量等于球b动能的增加量}
{相遇后的任意时刻，重力对球a做功功率和对球b做功功率相等}

%% answer: C
%% memo:
\memoanswer{
相遇时间为 $t$，则有 $\frac{h}{2}=\frac{1}{2}gt^{2}$，$\frac{h}{2}=v_{0}t-\frac{1}{2}gt^{2}$，两式联立得 $t=\frac{h}{v_{0}}$，相遇时球 a 的速度为 $gt=g\frac{h}{v_{0}}$，球 b 的速度为 $v_{0}-gt=v_{0}-g\frac{h}{v_{0}}$，故两者速度不一定相等、也不能同时落地，相遇之后的速度不同，故重力的功率也不相同；根据动能定理，两球重力做功分别为 $\frac{mgh}{2}$、$-\frac{mgh}{2}$，故 C 正确。
}

\item
%% number: 19
%% paperName: 2011年山东理综第 19 题 4 分
%% typeId: 2
%% type: 多选题
%% chapter: 运动和力
%% point: 瞬时加速度
%% method: 无
%% score: 4
%% degree: 450
%% duplicateId: 0
%% body:
如图所示，将两相同的木块a、b置于粗糙的水平地面上，中间用一轻弹簧连接，两侧用细绳固定于墙壁。开始时a、b均静止。弹簧处于伸长状态，两细绳均有拉力，a所受摩擦力$F_{fa}\neq0$，b所受摩擦力$F_{fb}=0$，现将右侧细绳剪断，则剪断瞬间\xzanswer{AD}

\onepicture[w=5.60cm]{figs/19.png}

\fourchoices[answer=AD]
{$F_{fa}$大小不变}
{$F_{fa}$方向改变}
{$F_{fb}$仍然为零}
{$F_{fb}$方向向右}

%% answer: AD
%% memo:
\memoanswer{
两物块相同，由受力分析可知两物体受到弹簧拉力大小相等，方向相反，绳子对 b 的拉力等于弹簧对 b 的拉力，若 a 平衡且有摩擦力，则绳对 a 拉力大小等于 b 受到的绳子拉力大小，此摩擦力小于最大静摩擦力，故当剪断右侧细绳后，a 受力情况不变，b 受到的摩擦力与弹簧拉力平衡，大小与 a 受到的摩擦力相等。

\twopicture[wA=6.00cm, wB=7.40cm]{figs/19_an1.png}{figs/19_an2.png}
}

\item
%% number: 20
%% paperName: 2011年山东理综第 20 题 4 分
%% typeId: 2
%% type: 多选题
%% chapter: 交流电
%% point: 变压器
%% method: 无
%% score: 4
%% degree: 450
%% duplicateId: 0
%% body:
为保证用户电压稳定在$220\UV$，变电所需适时进行调压，图甲为调压变压器示意图。保持输入电压$u_{1}$不变，当滑动接头P上下移动时可改变输出电压。某次检测得到用户电压$u_{2}$随时间$t$变化的曲线如图乙所示。以下正确的是\xzanswer{BD}

\twopicture[numA=甲, numB=乙, wA=4.50cm, wB=5.25cm]{figs/20a.png}{figs/20b.png}

\fourchoices[answer=BD]
{$u_{2}=190\sqrt{2}\sin(50\pi t)\UV$}
{$u_{2}=190\sqrt{2}\sin(100\pi t)\UV$}
{为使用户电压稳定在$220\UV$，应将P适当下移}
{为使用户电压稳定在$220\UV$，应将P适当上移}

%% answer: BD
%% memo:
\memoanswer{
根据图像知周期为 $T=2\times10^{-2}\Us$，电压的最大值为 $190\sqrt{2}\UV$，故用户得到电压的瞬时值为 $u_{2}=190\sqrt{2}\sin(100\pi t)\UV$，B 正确；用户获得电压为交流电的有效值，此时有效值为 $190\UV$，根据变压比 $\frac{n_{1}}{n_{2}}=\frac{U_{1}}{U_{2}}$，增大电压则需要减少原线圈的匝数 P 要适当上移，D 正确。
}

\item
%% number: 21
%% paperName: 2011年山东理综第 21 题 4 分
%% typeId: 1
%% type: 单选题
%% chapter: 电场
%% point: 电场线
%% method: 无
%% score: 4
%% degree: 450
%% duplicateId: 0
%% body:
如图所示，在两等量异种点电荷的电场中，MN为两电荷连线的中垂线，a、b、c三点所在直线平行于两电荷的连线，且a与c关于MN对称，b点位于MN上，d点位于两电荷的连线上。以下判断正确的是\xzanswer{BC}

\onepicture[w=3.20cm]{figs/21.png}

\fourchoices[answer=BC]
{b点场强大于d点场强}
{b点场强小于d点场强}
{a、b两点的电势差等于b、c两点间的电势差}
{试探电荷$+q$在$a$点的电势能小于在c点的电势能}

%% answer: BC
%% memo:
\memoanswer{
根据等量异种电荷的电场线分布可知 b 点场强小于 d 点场强，B 正确，A 错误；由对称性可知 a、b 两点的电势差等于 b、c 两点间的电势差，C 正确；MN 左侧电势大于零，而右侧小于零，所以试探电荷 $+q$ 在 a 点的电势能大于在 c 点的电势能，D 错误。
}

\item
%% number: 22
%% paperName: 2011年山东理综第 22 题 4 分
%% typeId: 2
%% type: 多选题
%% chapter: 电磁感应
%% point: 双杆问题
%% method: 无
%% score: 4
%% degree: 450
%% duplicateId: 0
%% body:
如图甲所示，两固定的竖直光滑金属导轨足够长且电阻不计。两质量、长度均相同的导体棒c、d，置于边界水平的匀强磁场上方同一高度$h$处。磁场宽为$3h$，方向与导轨平面垂直。先由静止释放c，c刚进入磁场即匀速运动，此时再由静止释放d，两导体棒与导轨始终保持良好接触。用$a_{c}$表示c的加速度，$E_{kd}$表示d的动能，$x_{c}$、$x_{d}$分别表示c、d相对释放点的位移。图乙中正确的是\xzanswer{BD}

\onepicture[w=3.00cm]{figs/22a.png}

\fourchoices[answer=BD, ispicture=true, h=2.40cm]
{figs/22b.png}
{figs/22c.png}
{figs/22d.png}
{figs/22e.png}

%% answer: BD
%% memo:
\memoanswer{
c 导体棒落入磁场之前做自由落体运动，加速度恒为 g，有 $h=\frac{1}{2}gt^{2}$，$v=gt$，c 棒进入磁场以速度 v 做匀速直线运动时，d 棒开始做自由落体运动，与 c 棒做自由落体运动的过程相同，此时 c 棒在磁场中做匀速直线运动的路程为 $h'=vt=gt^{2}=2h$，d 棒进入磁场而 c 还没有传出磁场的过程，无电磁感应，两导体棒仅受到重力作用，加速度均为 g，直到 c 棒穿出磁场，B 正确。c 棒穿出磁场，d 棒切割磁感线产生电动势，在回路中产生感应电流，因此时 d 棒速度大于 c 进入磁场时切割磁感线的速度，故电动势、电流、安培力都大于 c 刚进入磁场时的大小，d 棒减速，直到穿出磁场仅受重力，做匀加速运动，结合匀变速直线运动 $v^{2}-v_{0}^{2}=2gh$，可知加速过程动能与路程成正比，D 正确。
}

\item
%% number: 23
%% paperName: 2011年山东理综第 23 题 6 分
%% typeId: 4
%% type: 实验题
%% chapter: 电路
%% point: 电路实验
%% method: 无
%% score: 6
%% degree: 450
%% duplicateId: 0
%% body:
\begin{enumerate}
\item
某探究小组设计了“用一把尺子测定动摩擦因数”的实验方案。如图所示，将一个小球和一个滑块用细绳连接，跨在斜面上端。开始时小球和滑块均静止，剪断细绳后，小球自由下落，滑块沿斜面下滑，可先后听到小球落地和滑块撞击挡板的声音，保持小球和滑块释放的位置不变，调整挡板位置，重复以上操作，直到能同时听到小球落地和滑块撞击挡板的声音。用刻度尺测出小球下落的高度$H$、滑块释放点与挡板处的高度差$h$和沿斜面运动的位移$x$。（空气阻力对本实验的影响可以忽略）
\begin{steps}
\item
滑块沿斜面运动的加速度与重力加速度的比值为\rule{1.5cm}{0.4pt}。
\item
滑块与斜面间的动摩擦因数为\rule{2.0cm}{0.4pt}。
\item
（多选）以下能引起实验误差的是\rule{1.5cm}{0.4pt}。

（A）滑块的质量

（B）当地重力加速度的大小

（C）长度测量时的读数误差

（D）小球落地和滑块撞击挡板不同时
\end{steps}
\onepicture[w=3.40cm, align=right]{figs/23.png}
\item
某同学利用图甲所示电路，探究了电源在不同负载下的输出功率。

\onepicture[w=3.40cm, align=right]{figs/23a.png}

\begin{steps}
\item
所得实验数据如下表，请在给出的直角坐标系上画出$U-I$的图像。

\begin{center}
\begin{tabular}{|c|c|c|c|c|c|c|}
\hline
$U/\UV$ & 1.96 & 1.86 & 1.80 & 1.84 & 1.64 & 1.56 \\
\hline
$I/\UA$ & 0.05 & 0.15 & 0.25 & 0.35 & 0.45 & 0.55 \\
\hline
\end{tabular}
\end{center}

\onepicture[w=5.20cm, align=right]{figs/23b.png}

\item
根据所画$U-I$的图像，可求得电流$I=0.20\UA$时电源的输出功率为\rule{1.5cm}{0.4pt}$\UW$。（保留两位有效数字）
\item
实验完成后，该同学对实验方案进行了反思，认为按图甲电路进行实验操作的过程中存在安全隐患，并对电路重新设计。在图乙所示的电路中，你认为既能测出电源在不同负载下的输出功率，又能消除安全隐患的是\rule{1.5cm}{0.4pt}。（$R_{x}$阻值未知）

\onepicture[w=8.60cm, align=right]{figs/23c.png}
\end{steps}
\end{enumerate}
%% answer: 
\jdanswer{
\begin{enumerate}
\item
$\frac{x}{H}$；$\left(h-\frac{x^{2}}{H}\right)\frac{1}{\sqrt{x^{2}-h^{2}}}$；CD

\item
见解析图；$0.368\UW$；b、c

\end{enumerate}
}
%% memo:
\memoanswer{
（1）滑块沿斜面做初速度为零的匀加速直线运动，有 $x=\frac{1}{2}at^{2}$，小球做自由落体运动，有 $H=\frac{1}{2}gt^{2}$，所以 $\frac{a}{g}=\frac{x}{H}$。对滑块受力分析，根据牛顿第二定律得 $a=g\sin\theta-\mu g\cos\theta$，即 $\frac{x}{H}g=g\frac{h}{x}-\mu g\frac{\sqrt{x^{2}-h^{2}}}{x}$，解得 $\mu=\left(h-\frac{x^{2}}{H}\right)\frac{1}{\sqrt{x^{2}-h^{2}}}$。由上述分析知，C、D 能引起实验误差。

（2）$U-I$ 图像如图所示。

\onepicture[w=7.00cm]{figs/23_an.png}

由 $U-I$ 图像可知，当 $I=0.2\UA$ 时，$U=1.84\UV$。由 $P=UI$ 得：$P=1.84\times0.2\UW=0.368\UW$。按图甲电路进行实验操作的过程中存在安全隐患是：滑动变阻器的滑片移动到最右端时，相当于把电流表直接接在电源两极，会烧坏电流表。而图乙中的 a 仍然出现图甲的问题；b 中的 $R_{x}$ 在滑动变阻器的滑片移动时起分流作用保护电流表；c 中的 $R_{x}$ 在滑动变阻器的滑片移动到最右端时起限流作用保护电流表；d 中的 $R_{x}$ 在滑动变阻器的滑片移动到最右端时，$R_{x}$ 与 r 串联相当于电源内阻，等于直接把电流表接在电源两极。b、c 正确。
}

\item
%% number: 24
%% paperName: 2011年山东理综第 24 题 15 分
%% typeId: 6
%% type: 计算题
%% chapter: 运动和力
%% point: 连接体
%% method: 无
%% score: 15
%% degree: 450
%% duplicateId: 0
%% body:
如图所示，在高出水平地面$h=1.8\Um$的光滑平台上放置一质量$M=2\Ukg$、由两种不同材料连接成一体的薄板A，其右段长度$l_{1}=0.2\Um$且表面光滑，左段表面粗糙。在A最右端放有可视为质点的物块B，其质量$m=1\Ukg$。B与A左段间动摩擦因数$\mu=0.4$。开始时二者均静止，现对A施加$F=20\UN$水平向右的恒力，待B脱离A（A尚未露出平台）后，将A取走。B离开平台后的落地点与平台右边缘的水平距离$x=1.2\Um$。（取$g=9.8\Umsq$）求：
\begin{enumerate}
\item
B离开平台时的速度$v_{B}$；
\item
B从开始运动到刚脱离A时，B运动的时间$t_{B}$和位移$x_{B}$；
\item
A左端的长度$l_{2}$。
\end{enumerate}
\onepicture[w=6.60cm, align=right]{figs/24.png}
%% answer: 
\jdanswer{
\begin{enumerate}
\item
$2\Ums$

\item
$0.5\Us$，$0.5\Um$

\item
$1.5\Um$

\end{enumerate}
}
%% memo:
\memoanswer{
（1）B 离开平台做平抛运动。竖直方向有 $h=\frac{1}{2}gt^{2}$ ①，水平方向有 $x=v_{B}t$ ②，由①②式解得 $v_{B}=x\cdot\sqrt{\frac{g}{2h}}$，代入数据求得 $v_{B}=2\Ums$ ③。

（2）设 B 的加速度为 $a_{B}$，由牛顿第二定律和运动学知识得 $\mu mg=ma_{B}$ ④，$v_{B}=a_{B}t$ ⑤，$x_{B}=\frac{1}{2}a_{B}t_{B}^{2}$ ⑥，联立③④⑤⑥式，代入数据解得 $t_{B}=0.5\Us$ ⑦，$x_{B}=0.5\Um$ ⑧。

（3）设 B 刚开始运动时 A 的速度为 $v_{1}$，由动能定理得 $Fl_{1}=\frac{1}{2}Mv_{1}^{2}$ ⑨，设 B 运动时 A 的加速度为 $a_{A}$，由牛顿第二定律和运动学知识有 $F-\mu mg=Ma_{A}$ ⑩，$l_{2}+x_{B}=v_{1}t_{B}+\frac{1}{2}a_{A}t_{B}^{2}$ \ccnum{⑪}，联立⑦⑧⑨⑩\ccnum{⑪}式，代入数据解得 $l_{2}=1.5\Um$ \ccnum{⑫}。
}

\item
%% number: 25
%% paperName: 2011年山东理综第 25 题 18 分
%% typeId: 6
%% type: 计算题
%% chapter: 磁场
%% point: 带电粒子在磁场中的运动
%% method: 无
%% score: 18
%% degree: 250
%% duplicateId: 0
%% body:
扭摆器是同步辐射装置中的插入件，能使粒子的运动轨迹发生扭摆。其简化模型如图：Ⅰ、Ⅱ两处的条形匀强磁场区边界竖直，相距为$L$，磁场方向相反且垂直纸面。一质量为$m$、电量为$-q$、重力不计的粒子，从靠近平行板电容器MN板处由静止释放，极板间电压为$U$，粒子经电场加速后平行于纸面射入Ⅰ区，射入时速度与水平方向夹角$\theta=30^{\circ}$。
\begin{enumerate}
\item
当Ⅰ区宽度$L_{1}=L$、磁感应强度大小$B_{1}=B_{0}$时，粒子从Ⅰ区右边界射出时速度与水平方向夹角也为$30^{\circ}$，求$B_{0}$及粒子在Ⅰ区运动的时间$t_{0}$；
\item
若Ⅱ区宽度$L_{2}=L_{1}=L$，磁感应强度大小$B_{2}=B_{1}=B_{0}$，求粒子在Ⅰ区的最高点与Ⅱ区的最低点之间的高度差$h$；
\item
若$L_{2}=L_{1}=L$、$B_{1}=B_{0}$，为使粒子能返回Ⅰ区，求$B_{2}$应满足的条件；
\item
若$B_{1}\neq B_{2}$，$L_{1}\neq L_{2}$，且已保证了粒子能从Ⅱ区右边界射出。为使粒子从Ⅱ区右边界射出的方向与从Ⅰ区左边界射出的方向总相同，求$B_{1}$、$B_{2}$、$L_{1}$、$L_{2}$之间应满足的关系式。
\end{enumerate}
\onepicture[w=8.00cm, align=right]{figs/25.png}
%% answer: 
\jdanswer{
\begin{enumerate}
\item
$B_{0}=\frac{1}{L}\sqrt{\frac{2mU}{q}}$，$t_{0}=\frac{\pi L}{3}\sqrt{\frac{m}{2qU}}$

\item
$\left(2-\frac{2\sqrt{3}}{3}\right)L$

\item
$B_{2}\geqslant\frac{3}{L}\sqrt{\frac{mU}{2q}}$

\item
$B_{1}L_{1}=B_{2}L_{2}$

\end{enumerate}
}
%% memo:
\memoanswer{
（1）如图 1 所示，设粒子射入磁场Ⅰ区的速度为 $v$，在磁场Ⅰ区做圆周运动半径为 $R_{1}$，由动能定理和牛顿第二定律得

$qU=\frac{1}{2}mv^{2}$ ①

$qvB_{0}=m\frac{v^{2}}{R_{1}}$ ②

由几何关系得 $R_{1}=L_{2}=L$ ③

联立①②③得 $B_{0}=\frac{1}{L}\sqrt{\frac{2mU}{q}}$ ④

\onepicture[w=7.00cm]{figs/25_an1.png}

设粒子在Ⅰ区做圆周运动周期为 $T$，运动时间为 $t$，$T=\frac{2\pi R_{1}}{v}$ ⑤，$\frac{t}{T}=\frac{2\theta}{360^{\circ}}$ ⑥，联立①③⑤⑥式解得 $t=\frac{\pi L}{3}\sqrt{\frac{m}{2qU}}$ ⑦。

（2）设粒子在磁场Ⅱ区做圆周运动半径为 $R_{2}$，由牛顿第二定律得

$qB_{2}R_{2}=m\frac{v^{2}}{R_{2}}$ ⑧

由几何知识得 $h=(R_{1}+R_{2})(1-\cos\theta)+L\tan\theta$ ⑨

联立②③⑧⑨式解得 $h=\left(2-\frac{2\sqrt{3}}{3}\right)L$ ⑩

（3）如图 2 所示，为使粒子能再次返回到Ⅰ区应满足 $R_{2}(1+\sin\theta)\leq L$ \ccnum{⑪}

联立①⑧\ccnum{⑪}式解得 $B_{2}\geqslant\frac{3}{L}\sqrt{\frac{mU}{2q}}$ \ccnum{⑫}

\onepicture[w=7.20cm]{figs/25_an2.png}

（4）如图 3（或图 4）所示，设粒子射出磁场Ⅰ区时速度与水平方向的夹角为 $\alpha$，由图 2 几何知识可得

$L_{1}=R_{1}(\sin\theta+\sin\alpha)$ 或 $L_{1}=R_{1}(\sin\theta-\sin\alpha)$ \ccnum{⑬}

$L_{2}=R_{2}(\sin\theta+\sin\alpha)$ 或 $L_{2}=R_{2}(\sin\theta-\sin\alpha)$ \ccnum{⑭}

联立②⑧\ccnum{⑬}\ccnum{⑭}式解得 $B_{1}L_{1}=B_{2}L_{2}$ \ccnum{⑮}

\twopicture[wA=7.00cm, wB=6.80cm]{figs/25_an3.png}{figs/25_an4.png}
}

\item
%% number: 36
%% paperName: 2011年山东理综第 36 题 4 分
%% typeId: 6
%% type: 计算题
%% chapter: 气体性质
%% point: 查理定律
%% method: 无
%% score: 4
%% degree: 650
%% duplicateId: 0
%% body:
\begin{enumerate}
\item
人类对物质属性的认识是从宏观到微观不断深入的过程。以下说法正确的是\xzanswer{AD}
\fourchoices[answer=AD]
{液体的分子势能与体积有关}
{晶体的物理性质都是各向异性的}
{温度升高，每个分子的动能都增大}
{露珠呈球状是由于液体表面张力的作用}

\item
气体温度计结构如图所示。玻璃测温泡A内充有理想气体，通过细玻璃管B和水银压强计相连。开始时A处于冰水混合物中，左管C中水银面在O点处，右管D中水银面高出O点$h_{1}=14\Ucm$。后将A放入待测恒温槽中，上下移动D，使C中水银面仍在O点处，测得D中水银面高出O点$h_{2}=44\Ucm$。（已知外界大气压为1个标准大气压，1标准大气压相当于$76\UcmHg$）求恒温槽的温度。此过程A内气体内能\rule{1.5cm}{0.4pt}（填“增大”或“减小”），气体不对外做功，气体将\rule{1.5cm}{0.4pt}（填“吸热”或“放热”）。

\onepicture[w=2.80cm, align=right]{figs/36.png}
\end{enumerate}
%% answer: 
\jdanswer{
\begin{enumerate}
\item
AD

\item
$T_{2}=364\UK$；增大，吸热

\end{enumerate}
}
%% memo:
\memoanswer{
（1）物体体积变化时，分子间的距离将发生改变，分子势能随之改变，所以分子势能与体积有关，A 正确。晶体分为单晶体和多晶体，单晶体的物理性质各向异性，多晶体的物理性质各向同性，B 错误。温度是分子平均动能的标志，具有统计的意义，C 错误。液体表面的张力具有使液体表面收缩到最小的趋势，D 正确。

（2）①设恒温槽的温度为 $T_{2}$，由题意知 $T_{1}=273\UK$，A 内气体发生等容变化。由查理定律得 $\frac{P_{1}}{T_{1}}=\frac{P_{2}}{T_{2}}$ ①，$P_{1}=P_{0}+P_{h_{1}}$ ②，$P_{2}=P_{0}+P_{h_{2}}$ ③，联立①②③式解得 $T_{2}=364\UK$ ④。②理想气体的内能只由温度决定，A 气体的温度升高，所以内能增大。由热力学第一定律 $\Delta U=Q+W$ 知，气体不对外做功，气体将吸热。
}

\item
%% number: 37
%% paperName: 2011年山东理综第 37 题 4 分
%% typeId: 3
%% type: 填空题
%% chapter: 光学
%% point: 光的折射
%% method: 无
%% score: 4
%% degree: 650
%% duplicateId: 0
%% body:
\begin{enumerate}
\item
如图所示，一列简谐波沿$x$轴传播，实线为$t_{1}=0$时的波形图，此时P质点向$y$轴负方向运动，虚线为$t_{2}=0.01\Us$时的波形图。已知周期$T>0.01\Us$。波沿$x$轴\rule{1.5cm}{0.4pt}（填“正”或“负”）方向传播。求波速。

\onepicture[w=4.60cm, align=right]{figs/37a.png}

\item
如图所示，扇形AOB为透明柱状介质的横截面，圆心角$\angle$AOB$=60^{\circ}$。一束平行于角平分线OM的单色光由OA射入介质，经OA折射的光线恰平行于OB。求介质的折射率。折射光线中恰好射到M点的光线\rule{1.5cm}{0.4pt}（填“能”或“不能”）发生全反射。

\onepicture[w=4.00cm, align=right]{figs/37b.png}
\end{enumerate}
%% answer: 
\jdanswer{
\begin{enumerate}
\item
正；$100\Ums$

\item
$\sqrt{3}$；不能

\end{enumerate}
}
%% memo:
\memoanswer{
（1）$t_{1}=0$ 时，P 质点向 $y$ 轴负方向运动，由波的形成与传播知识可以判断波沿 $x$ 轴正向传播。由题意知 $\lambda=8\Um$，$t_{2}-t_{1}=\frac{1}{8}T$ ①，$v=\frac{\lambda}{T}$ ②，联立①②式代入数据求得 $v=100\Ums$。

（2）依题意作出光路图，如图所示。由几何知识知入射角 $i=60^{\circ}$，折射角 $\gamma=30^{\circ}$。根据折射定律 $n=\frac{\sin i}{\sin\gamma}$ 代入数据求得 $n=\sqrt{3}$。

\onepicture[w=6.00cm]{figs/37_an1.png}

$i'=30^{\circ}$，$\sin i'=\frac{1}{2}$，而 $\sin C=\frac{1}{n}=\frac{\sqrt{3}}{3}$。$i'<C$，所以不能发生全反射。

\onepicture[w=5.80cm]{figs/37_an2.png}
}

\item
%% number: 38
%% paperName: 2011年山东理综第 38 题 4 分
%% typeId: 6
%% type: 计算题
%% chapter: 运动和力
%% point: 动量守恒定律
%% method: 无
%% score: 4
%% degree: 650
%% duplicateId: 0
%% body:
\begin{enumerate}
\item
碘131核不稳定，会发生$\beta$衰变，其半衰变期为8天。碘131核的衰变方程：$\ce{^{131}_{53}I}\to$\rule{2.6cm}{0.4pt}（衰变后的元素用X表示）。经过\rule{1.5cm}{0.4pt}天有75\%的碘131核发生了衰变。
\item
如图所示，甲、乙两船的总质量（包括船、人和货物）分别为10$m$、12$m$，两船沿同一直线同一方向运动，速度分别为$2v_{0}$、$v_{0}$。为避免两船相撞，乙船上的人将一质量为$m$的货物沿水平方向抛向甲船，甲船上的人将货物接住，求抛出货物的最小速度。（不计水的阻力）

\onepicture[w=5.60cm, align=right]{figs/38.png}
\end{enumerate}
%% answer: 
\jdanswer{
\begin{enumerate}
\item
$\ce{^{131}_{54}X + ^{0}_{-1}e}$；16 天

\item
$4v_{0}$

\end{enumerate}
}
%% memo:
\memoanswer{
（1）核衰变遵守电荷量和质量数守恒，所以 $\ce{^{131}_{53}I -> ^{131}_{54}X + ^{0}_{-1}e}$。根据半衰期概念，得 $m_{\text{剩}}=m_{0}\cdot\left(\frac{1}{2}\right)^{\frac{t}{\tau}}$，即 $25\%m_{0}=m_{0}\cdot\left(\frac{1}{2}\right)^{\frac{t}{8}}$，解得 $t=16$ 天。

（2）设乙船上的人抛出货物的最小速度为 $v_{\min}$，抛出货物后的速度为 $v_{1}$，甲船上的人接到货物后速度为 $v_{2}$，由动量守恒定律得：$12m\times v_{0}=11m\times v_{1}-m\times v_{\min}$ ①，$10m\times 2v_{0}-m\times v_{\min}=11m\times v_{2}$ ②，为避免两船相撞应满足：$v_{1}=v_{2}$ ③，联立①②③式得：$v_{\min}=4v_{0}$ ④。
}

\end{enumerate}

\end{document}
